\section[Introdução]{Introdução}

Este capítulo apresenta uma análise reflexiva da minha trajetória na Agência Experimental de Engenharia de Software I (AGES I), realizada no segundo semestre de 2023. Durante esse período, participei do projeto \textbf{Plan Line}, uma aplicação personalizada de \textit{planning poker} desenvolvida para a empresa Triider.

Antes de iniciar a descrição do projeto, considero importante contextualizar o momento em que cheguei à AGES I. Os anos anteriores foram marcados por desafios pessoais e pela turbulência causada pela pandemia, o que me levou a refletir profundamente sobre meus objetivos profissionais. Após esse período, decidi ingressar no curso de Engenharia de Software, mesmo sem experiência prévia em desenvolvimento de software. Assim, a AGES I representou meu primeiro contato mais concreto com um projeto de software real, com equipe, stakeholders, prazos, tecnologias e responsabilidades compartilhadas.

Ao final de julho de 2023, após avaliar as oportunidades oferecidas pela AGES, escolhi participar do Plan Line. A escolha foi motivada pelo caráter técnico do projeto e pela possibilidade de aprender, na prática, conceitos fundamentais de desenvolvimento web e metodologias ágeis. O projeto foi orientado pelo professor Edson Moreno e teve como stakeholders Cassio, Letícia e Augusto, representantes da Triider.

A proposta do Plan Line era permitir que a equipe da Triider importasse as \textit{user stories} de uma sprint e realizasse a pontuação dessas histórias de acordo com as necessidades do projeto. A dinâmica esperada era semelhante à de ferramentas como Kahoot e Quizizz: os participantes votariam preferencialmente por dispositivos móveis, enquanto uma pessoa compartilharia a tela no computador para que todos acompanhassem as pontuações e resultados da sessão.

Essa experiência me proporcionou contato com uma aplicação real, com demandas de cliente, organização em sprints e necessidade de colaboração entre diferentes frentes técnicas. Além disso, permitiu explorar aspectos práticos da metodologia ágil, especialmente relacionados à estimativa de esforço, comunicação com stakeholders e adaptação contínua do produto às necessidades do usuário.
